\subsection*{13.4 Conditional Expectation Given a Discrete Random Variable}

Variable

Although conditional expectation can be defined in a very general setting, we often

want to compute it explicitly in concrete cases, such as when dealing with discrete

random variable. To do so, we need the following property of Lebesgue integral.

\begin{thmbox}{13.11}
\end{thmbox}

Suppose X is an integrable random variable defined on a probability space

(\Omega,\mathcal{F},P) and.(Ai)\infty

i=1 be a sequence of mutually disjoint sets in \mathcal{F}Then

\infty\sum

i=1

\int

Ai

XdP =

\int

\cupiAi

XdP.

\textit{Proof} For each integer m\geq 1 ,l e t Ym =X \summ

i=1 1Ai We obviously have

\vertYm\vert\leq\vert X\vert, and since E[\vertX\vert] is finite by assumption, we can apply the dominated

convergence theorem to obtain

\infty\sum

i=1

\int

Ai

XdP = limm\to\infty

m\sum

i=1

\int

\Omega

X1Ai dP=

\int

\Omega

limm\to\infty

m\sum

i=1

X1Ai dP=

\int

\cupiAi

XdP.

This proves the equality in Theorem 13.11. \hfill $\square$

With the result in the previous theorem, we can compute the conditional

expectation given a countable partition of the sample space.

\begin{thmbox}{13.12}
\end{thmbox}

Suppose that X is an integrable random variable defined on a probability

space.(\Omega,\mathcal{F},P) , andAi \in\mathcal{F},f o ri= 1 ,2,3,... , form a countable partition

of \Omega Let \mathcal{G}be the \sigma-algebra generated by .(Ai)\infty

i=1. Then, the conditional

expectationE[X\vert\mathcal{G}] is almost surely equal to .

\sum\infty

i=1 E[X\vertAi]1Ai .

\textit{Proof} Since the events Ai 's are mutually disjoint and their union is \Omega ,t h e

conditional expectation E[X\vert\mathcal{G}] is a function that is constant in each Ai Suppose

E[X\vert\mathcal{G}]()= c i for \in A i We can determine the value of ci from the defining

property of conditional expectation

ciP(Ai)=

\int

Ai

E[X\vert\mathcal{G}]dP =

\int

Ai

XdP.

Therefore,ci = 1

P(Ai)

\int

Ai XdP = E [X\vertAi]We define a \mathcal{G}-measurable random

variable

X()\coloneqq

\infty\sum

i=1

E[X\vertAi]1Ai ().

By linearity of expectation, for any finite union B= A i1 \cupAi2 \cup\cdot\cdot\cdot\cup Aim , where

i1,i2,...,i m are m distinct indices, we have

\int

B

XdP =

m\sum

k=1

\int

Aik

XdP =

m\sum

k=1

\int

Aik

XdP =

\int

B

XdP. (13.12)

For countably many indices, we can obtain the same equality by Theorem 13.11.

This completes the proof that X is a version of E[X\vert\mathcal{G}]\hfill $\square$

When the countable partition is induced by a discrete random variable, Theo-

rem 13.12 can be rephrased as follows.

\begin{thmbox}{13.13}
\end{thmbox}

Suppose X and Y are discrete random variables taking values in

{0,1,2,..., }Let the joint pmf of X and Y be pXY (x, y), and the marginal

distribution of Y be pY (y)= \sum

x\geq 0 pXY (x, y)When g(X) is integrable, we

have

E[g(X)\vertY= y ]=

\sum

x\geq 0

g(x)pX\vertY (x\verty),

where pX\vertY (x\verty) is defined as pXY (x,y)

pY (y) In particular , if X is integrable, the

conditional expectation of X given Y is

E[X\vertY= y ]=

\sum

x\geq 0

xpX\vertY (x\verty).

\textit{Proof} Let Ay be the event .{: Y() = y },f o r y= 0 ,1,2,3... F r o m

Theorem 13.12,f o r\in A y ,we have

E[g(X)\vertY]()= 1

P(Ay)

\int

Ay

g(x)dP = 1

pY (y)

\sum

x\geq 0

g(x)pXY (x, y).

Thus, we can write

E[g(X)\vertY= y ]=

\sum

x\geq 0

g(x) pXY (x, y)

pY (y) =

\sum

x\geq 0

g(x)pX\vertY (x\verty)

fory= 0 ,1,2,... \hfill $\square$
